

\section{Positive total scale: the recovered finite-pitchfork fit}

The September archive contains an explicit full-parent theorem for a
finite nonzero commuting edge, even though it does not contain the later
A5 terminal files under their recorded names. Its source hypotheses are
automatic on the ordinary and finite-pitchfork branch with
$d,C,A,D,\alpha$ units. This note proves that exact implication and
supplies explicit full-parent projective boundary forms. Thus that
subbundle of A5 is now mathematically audited. Weak-$\alpha$,
projective-pitchfork, and reciprocal terminal branches remain separate
obligations; there is no whole-A5 admission.

\begin{theorem}
Let the original normalized C1 or C2 source be bounded, nonnegative and
physical over a real closed coefficient extension of $\mathbb R$.
Assume $d,C,A,D,\alpha$ are positive bounded units. Retain its exact
canonical zero coefficients, in particular $\lambda=0$, and every
other nonnegative original coefficient $n,m,\beta$. Then its entire
parent is tensor-SOS at every independent projective testing pair and
every real or complex testing state.
\end{theorem}

The proof uses the admitted separated strict-Gram lemma, the freshly
audited noncommuting-residue theorem A6, and the preserved finite-edge
critical-center theorem whose proof is reviewed below. All source and
testing rates remain independent.

\section{Exact source fit of the recovered finite-edge theorem}

The archived notation is
$(d,\lambda,e_1,e_2,e_3,e_4,e_5,\alpha,\beta)
=(d,\lambda,n,m,C,A,D,\alpha,\beta)$.
The theorem requires $\lambda=0$ exactly, $\bar d>0$, $\bar m=0$,
a finite real $v_0\ne0$, and positivity conditions on the residue
determinant at $u=0$. Bars denote real residues in this section.
Its test state is the $W$ ray. It permits arbitrary independent
nonnegative opening rates of all the original coefficients.

At $\lambda=m=0$, the exact determinant polynomial is
\[
 Q(U,V)=U E(V)+U^2 B(V),\quad
 E(V)=q_{10}+q_{11}V+q_{12}V^2,\quad
 B(V)=q_{20}+q_{21}V+q_{22}V^2.
\]
Use the inherited positive quantities
\begin{align*}
 t_1&=d+2(n+D),&t_2&=d+2(\alpha+m+A),\\
 t_3&=d+2(C+D),&t_4&=d+2(\alpha+\beta+m),\\
 t_5&=d+2(\beta+n),&t_6&=d+2(\alpha+C+A),\\
 t_7&=d+2(\beta+n+m),&t_8&=d+2(C+A+D).
\end{align*}
All eight have positive bounded-unit residues. At the residue $m=0$,
the required coefficients are
\[
 q_{10}=4\alpha A t_7t_3>0,\quad
 q_{12}=4CDt_5t_2>0,\quad
 q_{21}=4(\beta n t_3t_2+CA t_1t_4)>0,
\]
and $q_{20},q_{22}\geq0$. Therefore $B(v_0^2)>0$ for every
$v_0\ne0$. Physicality implies $E(V)\geq0$ for all real $V>0$:
divide $Q(U,V)\geq0$ by $U>0$ and let the real variable $U$ tend to
zero. Consequently a zero $E(v_0^2)=0$ is a double zero,
$E(V)=q_{12}(V-v_0^2)^2$, and otherwise $E(v_0^2)>0$.

After normalizing the tested event
$r_2=vW-X+uvY-uZ$ to one, the residue state complement on
$(X,Y,Z)$ is
\[
 M_0(v)=\begin{pmatrix}
 2(d+2n+2\beta)(1+v^2)&0&0\\
 0&2d(1+v^2)+4(C+A)(1+v^2)+4D+4\alpha v^2
                                      &4v(D-\alpha)\\
 0&4v(D-\alpha)&4\alpha+4Dv^2
 \end{pmatrix}.
\]
Its lower-block determinant is positive because, after a positive factor
is removed, it is
\[
 \alpha d+2\alpha C+2\alpha A+2\alpha D(1+v^2)
                      +Dv^2(d+2C+2A)>0.
\]
Thus the specified state complement is positive definite. Every
hypothesis of the recovered finite-edge theorem is now verified
directly from the original A5 source, including arbitrary $n,\beta\geq0$.

\paragraph{What the recovered proof actually establishes.}
Its exact pure-testing normalization is
$W=(1+X-uvY+uZ)/v$, leaving the entire opened $m$ packet as $4m$.
The formal coefficient-side substitution
$X=u^2x,Y=ua,Z=ub,U=u^2$ gives the full identity
\[
 F=4m+u^2K(x,a,b,v,U),\qquad K_x=U E_x,
\]
with $\partial_xE_x|_{U=0}=4(d+2n+2\beta)(1+v^2)>0$.
The actual independent testing coordinates are never divided by $u$.
The minimized residue profile is a positive multiple of $E(v^2)$;
the full $(a,b,v)$ Hessian is strict at a double zero. Its next
coefficient is
\[
 \gamma=\frac{16(1+v_0^2)^2 B(v_0^2)}{v_0^2\det M_0(v_0)}>0.
\]
This is a real residue determinant/Schur calculation, not a mixed
inverse in the tensor ring.

The coefficient-field implicit-function equations are
$(E_x,K_a,K_b,K_v)=0$. With $\kappa$ the optimized value at $U=0$,
use $U_c=0$ if $\kappa\geq0$; if $\kappa<0$, also impose
$K+UK_U=0$. The added pivot is $2\gamma>0$ and produces $U_c>0$.
The remaining scalar $\mu=4m+U_cK_c$ is nonnegative by original
physicality at an actual coefficient-field realization.

The full normal expansion uses
\[
 N_x=X-u^2x_c,\quad N_y=Y-ua_c,\quad N_z=Z-ub_c,
 \quad N_v=u(v-v_c),\quad R=u^2-U_c.
\]
Every first state-normal coefficient is printed in the recovered proof:
$u^2E_x,uK_a,uK_b$. Monic division at the coefficient center expresses
them in the five roots above, with an extra $u$ or $u^2$ on every
coupling to $R$. The complete state remainder is exactly quadratic.
The residue Gram consists of the positive state-$X$ mass, the positive
$(a,b,v)$ Hessian, and the pivot $\gamma$; all couplings to $R$ vanish
in residue. If $\kappa\geq0$, replace $R$ by $u^2$ and retain
$4m+\kappa u^2$. Otherwise retain $\mu$. Both branches restore the
original event normalization. This is a full-parent certificate.
When $E(v_0^2)>0$ the four ordinary normals $(X,Y,Z,u)$ already
have strict residue, and the same whole $4m$ packet is retained.

The complete proof is the recovered
\path{POSITIVE_D_FINITE_EDGE_TERMINAL.md}; its exact full-parent
checker was freshly replayed: 15 identities, 30 scalar entries pass.
Its local conclusion is used only on the finite nonzero edge stated
above. The next section supplies the previously excluded endpoints.

\section{Explicit all-projective dispatch}

Noncommuting residue tests use A6. It remains that at least one angle
has residue zero or infinity. For finite $u_0\ne0$ and $v_0=0$,
the real residue parent is a sum of nonnegative event terms and the
core $2d[(uW-Y)^2+(X+uZ)^2]$. The unit $A,\alpha,C$ terms
force $Y=Z=W=0$, and the core then forces $X=0$.
Hence the residue state matrix is positive definite.

For $u_0=0$ and finite $v_0\ne0$, the unit $C,A,D,\alpha$ terms
and $d>0$ leave only the $W$ ray if $m$ is weak, while unit $m$
removes that ray too. In the weak case, positive residue slack is
immediate; at the $W$ ray, the previous section applies for either
$E(v_0^2)>0$ or $E(v_0^2)=0$.

For $u=\infty$ and finite nonzero $v$, set $p=1/u$ and homogenize
as $F_p=p^2F(1/p,v)$. Its real residue has core
\[
 2d[(W+vX)^2+(vY-Z)^2]
\]
and positive $4C(1+v^2)W^2+4A(1+v^2)Y^2$ packets. These
force every state coordinate to vanish at a zero, so the matrix is
positive definite. For $v=\infty$ and finite nonzero $u$, the
homogenized residue has core $2d(1+u^2)(X^2+Y^2)$ together with
$4C(uW-Y)^2+4D(uX-Z)^2$. Again it is positive definite.
These are exact residue identities and use no density argument.

It remains to cover the four double corners. The following displayed
quadratic forms are complete, with coefficient letters interpreted in
real residue. All coefficients that are merely nonnegative remain so.
At a unit state coordinate, normalize it in the testing field.

At $(u,v)=(0,0)$ the only possible zero state is the $W$ ray.
Put $W=1$ and retain the entire $m$ packet. The remainder lies
exactly in $(u,X,Y,Z)^2$, and its residue quadratic is
\[
 2(d+2n+2\beta)X^2+(4C+4D+2d)(u-Y)^2+4AY^2+4\alpha Z^2.
\]
It is strict. The testing $v$ is bounded infinitesimal in its
coefficients and is never divided out.

At $(u,v)=(\infty,0)$ the only possible zero state is the $X$ ray
when $n$ is weak; if $n$ is a unit there is no zero. Put $X=1$,
retain the entire $n$ packet in $F_p$, and use roots $(p,v,W,Y,Z)$.
The whole remainder belongs to their square ideal and its residue is
\[
 4CW^2+4AY^2+(4D+2d)(W+v)^2
                 +4\alpha Z^2+[4(m+\beta)+2d](p+Z)^2\succ0.
\]

At $(u,v)=(0,\infty)$ the only possible zero state is the $W$ ray
when $m$ is weak. Put $q=1/v$, $W=1$, and retain the entire $m$
packet in $F_q=q^2F(u,1/q)$. The whole remainder is in
$(u,X,Y,Z)^2$ and its residue quadratic is
\[
 (4n+4\beta+2d)X^2+4C(u-Y)^2+[4(A+\alpha)+2d]Y^2+4DZ^2\succ0.
\]
The testing $q$ remains a bounded infinitesimal coefficient.

At $(u,v)=(\infty,\infty)$ the only possible zero state is the $Z$
ray when $\beta$ is weak. Put $Z=1$ and retain the entire $\beta$
packet in $F_{pq}=p^2q^2F(1/p,1/q)$. The complete remainder lies
in $(p,q,W,X,Y)^2$ and its residue quadratic is
\[
 4CW^2+4AY^2+4D(X-p)^2+(2d+4n)X^2
                            +[2d+4(m+\alpha)](Y-q)^2\succ0.
\]

For each corner, allocate two normal factors from every monomial to
the Gram roots; every remaining factor is bounded. Terms above
quadratic order have infinitesimal coefficients. The displayed full
quadratic residues are strict, so the separated strict-Gram lemma
applies. Restore each retained packet. Homogeneous angle factors and
state normalizations are pure testing squares. The identities include
literal zero and infinity and every original nonnegative opening rate.
For nonzero state residues away from the listed rays, positivity gives
the scalar strict-residue exit. The zero state is immediate and complex
states split into real and imaginary evaluations. This proves the theorem.

\section*{Scope, checks, and the remaining A5 imports}

The new checker \path{check_finite_pitchfork_fit.py} verifies 11 exact
identities: four full-remainder square-ideal memberships, their four
complete quadratic forms, and the three finite/infinite edge identities.
The recovered finite-edge checker separately verifies the exact
normalized full parent, all first state normals, and complete hard tail.
The general signs are proved above and in the recovered critical-center
argument; the checker is corroboration rather than a positivity premise.

This proves the ordinary and finite-pitchfork subbundle with
$d,C,A,D,\alpha$ units. In the C2 partition its source-reflection
image covers the single weak-$D$ branch. In the C1 ordinary branch a
unit $\beta$ may be exchanged with $D$ by the exact stabilizer when
$m,C,A,\alpha$ are units. These are the previously audited exact
full-parent symmetry transports, with their literal zeros retained.

The remaining A5 imports include the weak-$\alpha$ all-tests terminals and
their stabilizer images, the C1 projective pitchfork with $D,\beta$
weak, and the reciprocal source branches with $m,D$ weak. None satisfies
the required unit-$D$ premise by mere approximation. A separate note
\path{weak_alpha_unit_n_fit.tex} verifies the finite-ratio corner's
positive-endpoint and simple-complement hypotheses only on the C2
subdomain where $d,n,m,C,A,D$ are all units and $\alpha$ is weak.
It does not supply the rest of weak-$\alpha$.
The strict-$\alpha$-mass theorem requires $C,A\gg d$, which is
incompatible with bounded normalized coefficients and unit $d$.

Exact archive identity, selected member hashes, fresh replay receipts,
and the scoped G1--G9 review are recorded in the adjacent recovery
manifest and \path{RECOVERY_AND_FIT.txt}. The newer R19 frontier
continues to override old global-status assertions.
